% Standalone LaTeX source. Bibliography, tables and plots are embedded.
% Compile twice with pdflatex. No external bibliography or figures are required.
\documentclass[11pt,a4paper]{article}
\usepackage[T1]{fontenc}
\usepackage[utf8]{inputenc}
\usepackage{lmodern}
\usepackage[margin=24mm,headheight=14pt]{geometry}
\usepackage{amsmath,amssymb,amsthm,bm,mathtools}
\usepackage{booktabs,array,tabularx,longtable}
\usepackage{graphicx,xcolor,microtype,enumitem,listings}
\usepackage{tikz,pgfplots}
\pgfplotsset{compat=1.18}
\usepackage{fancyhdr,placeins}
\usepackage[colorlinks=true,linkcolor=blue!45!black,citecolor=blue!45!black,
urlcolor=blue!45!black,pdftitle={WJ-EARSM budget-cutoff analytical wall function}]{hyperref}
\usepackage[nameinlink,noabbrev]{cleveref}
\numberwithin{equation}{section}
\newtheorem{proposition}{Proposition}[section]
\newtheorem{remark}[proposition]{Remark}
\newcommand{\dd}{\mathrm d}
\newcommand{\eps}{\varepsilon}
\newcommand{\ust}{u_*}
\newcommand{\qe}{q_e}
\newcommand{\Ce}{C_*}
\newcommand{\ke}{\kappa_e}
\newcommand{\Ne}{N_e}
\newcommand{\tr}{\operatorname{tr}}
\newcommand{\diag}{\operatorname{diag}}
\newcommand{\I}{\mathsf I}
\newcommand{\R}{\mathsf R}
\newcommand{\units}[1]{\,\mathrm{#1}}
\newcommand{\pos}[1]{(#1)_+}
\setlength{\parindent}{0pt}
\setlength{\parskip}{4pt}
\setlength{\emergencystretch}{3em}
\setlist[itemize]{topsep=4pt,itemsep=2pt,parsep=0pt}
\setlist[enumerate]{topsep=4pt,itemsep=3pt,parsep=0pt}
\pagestyle{fancy}
\fancyhf{}
\fancyhead[L]{\small WJ-EARSM neutral ASL and analytical wall functions}
\fancyhead[R]{\small Standalone appendix}
\fancyfoot[C]{\thepage}
\lstdefinestyle{py}{language=Python,basicstyle=\ttfamily\scriptsize,
keywordstyle=\color{blue!55!black},commentstyle=\color{black!55},
stringstyle=\color{green!35!black},showstringspaces=false,
breaklines=true,columns=fullflexible,keepspaces=true,
frame=single,rulecolor=\color{black!20},framerule=0.3pt}

\title{\bfseries A budget-determined analytical wall function\\
for the high-Reynolds-number WJ-EARSM\\[4pt]
\large Standalone appendix: calibrated coefficients, nonconstant inner
kinetic energy and conservative coupling}
\author{Standalone research appendix}
\date{10 November 2026}
\begin{document}
\maketitle
\begin{abstract}
An auxiliary wall layer is coupled to the exact neutral surface-layer
reference of the Wallin--Johansson explicit algebraic Reynolds-stress
model. The outer system uses its declared high-Reynolds-number equations;
the complete analytical wall function retains positive molecular
viscosity. Matching the dissipation coefficient and viscosity slope fixes
the auxiliary coefficients. An integrated reference-production budget
then determines a unique admissible dissipation cutoff. Ordinary
diffusion of a nonconstant inner kinetic-energy profile balances the
local source defect, with zero flux at both matching endpoints.
Closed formulas, a sufficient positivity and covariance-realizability
proof, full traction and pressure matching, and a conservative curved
extension are derived. A quantitative example and a separately named
dynamic-positivity option distinguish stationary reference properties
from general-flow guarantees.
\end{abstract}
\clearpage
{\small\setlength{\parskip}{0pt}\tableofcontents}
\newpage
\appendix
\section{A budget-cutoff AWF with WJ-compatible coefficients and a nonconstant inner kinetic-energy profile}
\label{ba:section}

This appendix specifies a wall-layer variant that preserves the neutral
Wallin--Johansson (WJ) equilibrium outside the modeled layer while allowing
the transported kinetic energy to vary inside it. The auxiliary viscosity
is calibrated to the outer equilibrium, a global production budget determines
the dissipation cutoff, and ordinary diffusion of the inner kinetic energy
balances the remaining local production--dissipation difference. No prescribed
redistribution is added to the inner kinetic-energy equation.

The outer constitutive law is the high-Reynolds-number WJ model
\cite{WJ2000}, coupled to the scalar transport convention used by
Baungaard et al.\ \cite{Baungaard2022}. The auxiliary viscosity and empirical
roughness correlation originate in the rough-wall analytical wall function
of Suga, Craft, and Iacovides \cite{SCI2006}. The coefficient calibration,
budget cutoff, retained inner transport, and conservative coupling below
specify the proposed combination. Every dimensional AWF quantity retains
the strictly positive molecular viscosity $\nu$.

\subsection{Scope, endpoint sampling, and coefficient calibration}
\label{ba:coefficients}

The outer analytical reference uses the declared high-Re equations: molecular
momentum and scalar diffusion are omitted there. Within the modeled layer,
$\nu>0$ is retained in momentum, turbulence diffusion, production, roughness
scaling, and all auxiliary integrals. Consequently, the result is an exact
reference for this particular coupled system; it is not a claim that the
same profiles solve outer equations with finite molecular diffusion.

Let $z=0$ be the effective numerical ground and $z=H>0$ the matching
surface. Subscripts $w$ and $o$ identify the modeled wall layer and outer
region. Select endpoint sampling and endpoint tail matching:
\begin{equation}
 \boxed{z_K=z_m=H,\qquad K=k_w(H)=k_o(H)>0.}
 \label{ba:sampling}
\end{equation}
The transported inner turbulence variable is $k_w$. Inner dissipation is
a dependent positive sink, not an independently transported field. In
particular, there is no inner dissipation equation to which a second
interface condition could be assigned.

The selected outer coefficients are
\begin{equation}
 c_1=1.8,\quad c_2=\frac59,\quad
 C_{\epsilon1}=1.44,\quad C_{\epsilon2}=1.82,\quad
 \sigma_k=1,\quad \sigma_\epsilon=1.3,\quad C_d=C_*.
 \label{ba:outer-coefficients}
\end{equation}
Here $C_d$ is the constant coefficient in the scalar diffusivity
$\nu_d=C_dk^2/\epsilon$; it is distinct from the generally variable
momentum coefficient $C_\mu^{\mathrm{eff}}$.

\paragraph{Derivation of the outer constant-stress state.}
For general three-dimensional mean flow, the selected constitutive law is
the full three-dimensional explicit tensor of
\cite[Eqs.~(15)--(18)]{Baungaard2022}, using the two-dimensional cubic
solution for $N$ as a surrogate. It is not the exact implicit
three-dimensional ARSM solution. In planar mean flow, including the
simple shear considered here, it reduces exactly to the WJ planar law
\cite{WJ2000,Baungaard2022}. The derivation below therefore involves no
three-dimensional-root approximation.

Set $\boldsymbol U=U(z)\boldsymbol e_x$, $G=U'(z)>0$, and
$q=kG/\epsilon$. With the normalization
\begin{equation}
 \mathsf a=R/k-\tfrac23I,\qquad
 S=\frac{k}{2\epsilon}
       [\nabla\boldsymbol U+(\nabla\boldsymbol U)^T],\qquad
 \Omega=\frac{k}{2\epsilon}
       [\nabla\boldsymbol U-(\nabla\boldsymbol U)^T],
 \label{ba:outer-normalization}
\end{equation}
the only nonzero components are
$S_{xz}=S_{zx}=q/2$ and $\Omega_{xz}=-\Omega_{zx}=q/2$.
Consequently $\operatorname{tr}S^2=q^2/2$,
$\operatorname{tr}\Omega^2=-q^2/2$, and
$S\Omega-\Omega S=(q^2/2)\operatorname{diag}(-1,0,1)$.
For $c_2=5/9$, put $c'=\tfrac94(c_1-1)$. The exact planar law is
\begin{equation}
 \mathsf a=\beta_1S+\beta_4(S\Omega-\Omega S),\qquad
 (\beta_1,\beta_4)=-\frac6{5(N^2+q^2)}(N,1),
 \label{ba:outer-planar-stress}
\end{equation}
where $N$ is the physical root, $N\geq c'>0$, of
\begin{equation}
 \boxed{N^3-c'N^2-\frac7{20}q^2N-c'q^2=0.}
 \label{ba:outer-cubic}
\end{equation}
Hence
\begin{equation}
 \begin{gathered}
 C(q)=\frac{3N}{5(N^2+q^2)},\qquad
 a_{xz}=-C(q)q,\qquad
 a_{xx}=-a_{zz}=\frac{C(q)q^2}{N},\qquad a_{yy}=0,\\
 P=-R_{xz}G,\qquad \frac P\epsilon=C(q)q^2,\qquad
 N=c'+\frac94\frac P\epsilon.
 \end{gathered}
 \label{ba:outer-shear-identities}
\end{equation}
The last identity follows by rewriting the cubic as
$(N-c')(N^2+q^2)=27Nq^2/20$. It is exact for this planar reduction;
production in a general three-dimensional implementation must instead
be contracted from the full explicit stress.

The high-Re scalar equations are
\begin{align}
 \frac{Dk}{Dt}
 &=P-\epsilon+\nabla\cdot
       \left(\frac{\nu_d}{\sigma_k}\nabla k\right),
 \label{ba:outer-k-equation}\\
 \frac{D\epsilon}{Dt}
 &=(C_{\epsilon1}P-C_{\epsilon2}\epsilon)\frac{\epsilon}{k}
   +\nabla\cdot\left(\frac{\nu_d}{\sigma_\epsilon}\nabla\epsilon\right),
 \qquad \nu_d=C_d\frac{k^2}{\epsilon}.
 \label{ba:outer-epsilon-equation}
\end{align}
For a stationary, horizontally homogeneous state with $k=K$ constant,
\eqref{ba:outer-k-equation} requires $P=\epsilon$. Therefore
\begin{equation}
 N_e=\frac94c_1,\qquad
 C(q_e)q_e^2=1,\qquad
 q_e^2=\frac{N_e^2}{(3/5)N_e-1},\qquad
 C_*:=q_e^{-2}.
 \label{ba:outer-equilibrium-root}
\end{equation}
The zero-streamwise-pressure-gradient momentum equation requires
$\mathrm dR_{xz}/\mathrm dz=0$. Choose its constant as
$R_{xz}=-u_*^2$. Because $a_{xz,e}=-1/q_e$, this fixes
$K=q_eu_*^2$; production then fixes $\epsilon=u_*^2G$.
Thus a logarithmic velocity entails
\begin{equation}
 U(z)=\frac{u_*}{\kappa}\ln\frac z{z_0},\qquad
 k=K=q_eu_*^2,\qquad
 \epsilon(z)=\frac{u_*^3}{\kappa z},\qquad
 \nu_{t,o}(z)=\kappa u_*z.
 \label{ba:outer-log-family}
\end{equation}
The normal stresses are constant as well, since
$a_{xx}=1/N_e$, $a_{yy}=0$, and $a_{zz}=-1/N_e$.
Consequently the reduced outer pressure is constant, or mechanical
pressure is hydrostatic if gravity is retained. The outer differential
equations leave $z_0$ free; the wall anchor below determines it.

The remaining condition comes from dissipation transport, not from the
stress law alone. On \eqref{ba:outer-log-family},
\begin{equation}
 \nu_d=\frac{C_d}{C_*}\kappa u_*z,\qquad
 \frac{\mathrm d}{\mathrm dz}
       \left(\frac{\nu_d}{\sigma_\epsilon}\epsilon'\right)
 =\frac{C_d}{C_*}\frac{u_*^4}{\sigma_\epsilon z^2}.
 \label{ba:outer-epsilon-diffusion}
\end{equation}
Substituting into \eqref{ba:outer-epsilon-equation} gives the residual
\begin{equation}
 \mathcal R_{\epsilon,0}
 =\left[C_{\epsilon1}-C_{\epsilon2}
       +\frac{C_d\kappa^2}{\sigma_\epsilon C_*^{3/2}}\right]
       \frac{\epsilon^2}{K}.
 \label{ba:outer-epsilon-residual}
\end{equation}
It vanishes exactly when
\begin{equation}
 \boxed{\kappa^2
 =\sigma_\epsilon(C_{\epsilon2}-C_{\epsilon1})
            \frac{C_*^{3/2}}{C_d}.}
 \label{ba:outer-kappa-condition}
\end{equation}
For the selected $C_d=C_*$, these derivations give
\begin{equation}
 \begin{gathered}
 N_e=\frac94c_1=\frac{81}{20},\qquad
 C_* =\frac{(3/5)N_e-1}{N_e^2}
      =\frac{572}{6561},\qquad q_e=C_*^{-1/2},\\
 \kappa_e^2=(C_{\epsilon2}-C_{\epsilon1})
                  \sigma_\epsilon\sqrt{C_*},\\
 q_e=3.386780140528718,\qquad
 C_*=0.0871818320377991,\qquad
 \kappa_e=0.3819178635467331.
 \end{gathered}
 \label{ba:outer-equilibrium}
\end{equation}
These relations, rather than rounded displays of $C_*$ or $\kappa_e$,
define the reference and recover the neutral-ASL compatibility relation
in \cite[Eqs.~(22)--(23)]{Baungaard2022}.

At each positive sample $K$, define reference scales
\begin{equation}
 u_*(K)=\sqrt{\frac{K}{q_e}},\qquad
 t_0(K)=\frac{K}{q_e},\qquad
 \mathcal A(K)=\frac{u_*(K)^3}{\kappa_e}.
 \label{ba:reference-scales}
\end{equation}
These are counterfactual equilibrium scales evaluated from the current
sample. In a forced, developing, or separated flow, $t_0(K)$ is not
imposed as the actual wall traction. Arguments $K$ are suppressed when
all quantities use the same sample. A fixed inlet reference is recovered
by fixing $K$ throughout the flat reference calculation.

Calibrate the dissipation tail and the asymptotic auxiliary-viscosity
slope to the WJ branch:
\begin{equation}
 \boxed{\begin{aligned}
 c_\ell^w&=\kappa_e q_e^{3/2}=2.380402945263298,\\
 c_\mu^w&=C_*=q_e^{-2}=0.0871818320377991,\\
 \alpha&=c_\mu^wc_\ell^w
       =\frac{\kappa_e}{\sqrt{q_e}}=0.2075278897562272.
 \end{aligned}}
 \label{ba:wall-coefficients}
\end{equation}
In particular, if $a:=\alpha\sqrt K$,
\begin{equation}
 \boxed{a=\kappa_e\sqrt{K/q_e}=\kappa_e u_*,
 \qquad
 \frac{K^{3/2}}{c_\ell^w}=\mathcal A=\frac{t_0^2}{a}.}
 \label{ba:calibration-identities}
\end{equation}
Matching $c_\ell^w$ supplies the outer dissipation coefficient. Matching
$c_\mu^w$ additionally supplies the outer viscosity slope; finite-height
traction continuity alone would not require equal slopes.

Let $h\geq0$ be equivalent sand-grain roughness, with $h<H$ for the
wall-function geometry. Retain the empirical intercept correlation
\cite[Eqs.~(24)--(25)]{SCI2006} without claiming a new roughness calibration:
\begin{equation}
 \begin{aligned}
 h^*&=\frac{h\sqrt K}{\nu},\qquad
 b=10.7\left[1-\left(\frac{h^*}{70}\right)^m\right],\\
 m&=\max\left\{
 0.5-0.4\left(\frac{h^*}{70}\right)^{0.7},\,
 1-0.79\left(\frac{h^*}{70}\right)^{-0.28}\right\}
 \qquad(h^*>0).
 \end{aligned}
 \label{ba:roughness-correlation}
\end{equation}
For $h=0$, set $b=10.7$ directly. The auxiliary coefficients are
\begin{equation}
 \boxed{\begin{aligned}
 z_v&=\frac{\nu b}{\sqrt K},\qquad
 \nu_{t,A}(z)=a\max(z-z_v,0),\\
 D_A(z)&=\nu+\nu_{t,A}(z),\qquad
 D_{k,w}(z)=\nu+\frac{\nu_{t,A}(z)}{\sigma_k}.
 \end{aligned}}
 \label{ba:viscosities}
\end{equation}
Thus $D_{k,w}=D_A$ for the selected $\sigma_k=1$. A negative $z_v$ is a
virtual intercept of the viscosity line, not a displacement of the ground.
All coefficients in \eqref{ba:viscosities} use the sampled $K$, while
the transported field $k_w(z)$ is allowed to vary. Replacing $K$ by
local $k_w(z)$ in these coefficient laws would define a different,
nonlinearly coupled model.

\subsection{Momentum, production, and the outer velocity anchor}
\label{ba:momentum}

Put $L_h(z)=\max(z-h,0)$ and evaluate the following integrals with the
calibrated, molecular-viscosity-retaining coefficient $D_A$:
\begin{equation}
 J_0(z)=\int_0^z\frac{\mathrm ds}{D_A(s)},\qquad
 J_1(z)=\int_0^z\frac{L_h(s)}{D_A(s)}\,\mathrm ds,\qquad
 J_2(z)=\int_0^z\frac{\mathrm ds}{D_A(s)^2}.
 \label{ba:integrals}
\end{equation}
For the original flat, quasi-steady AWF reduction, $F$ denotes the sampled
tangential convection and mechanical-pressure forcing, in acceleration
units. The rough-wall forcing is suppressed below $h$:
\begin{equation}
 (D_AU_A')'=F\mathbf1_{z>h},\qquad U_A(0)=0,\qquad U_A(H)=U_H.
 \label{ba:momentum-ode}
\end{equation}
Integration gives
\begin{equation}
 \boxed{\begin{aligned}
 t_A&=\frac{U_H-FJ_1(H)}{J_0(H)},\\
 U_A(z)&=t_AJ_0(z)+FJ_1(z),\qquad
 U_A'(z)=\frac{t_A+FL_h(z)}{D_A(z)},\\
 P_A(z)&=\nu_{t,A}(z)
          \left[\frac{t_A+FL_h(z)}{D_A(z)}\right]^2.
 \end{aligned}}
 \label{ba:forced-awf}
\end{equation}
For two tangential velocity components the first two formulas apply
componentwise, and production uses the squared vector norm. Negative
traction is retained when the actual forcing and matching velocity yield
reversal. The mathematical form of the pressure-gradient response is
unchanged by the coefficient calibration.

At the zero-forcing reference, impose the velocity anchor
\begin{equation}
 U_{H,0}=t_0J_0(H),\qquad
 U_{w,0}(z)=t_0J_0(z),\qquad
 p_{A,0}(z)=t_0^2\frac{\nu_{t,A}(z)}{D_A(z)^2}.
 \label{ba:reference-production}
\end{equation}
The lowercase symbol $p_{A,0}$ denotes reference production, not pressure.
Its integrated value is
\begin{equation}
 \boxed{\mathcal P_0=\int_0^H p_{A,0}(z)\,\mathrm dz
 =t_0^2[J_0(H)-\nu J_2(H)].}
 \label{ba:production-budget}
\end{equation}
The term $\nu J_2$ subtracts mean molecular dissipation from total
mean shear work. Omitting it would change the turbulent-energy budget.

For $z\geq H$, take
\begin{equation}
 \boxed{\begin{aligned}
 U_{o,0}(z)&=t_0J_0(H)+\frac{u_*}{\kappa_e}\ln\frac zH,\\
 k_{o,0}(z)&=K,\qquad
 \epsilon_{o,0}(z)=\frac{\mathcal A}{z},\qquad
 \nu_{t,o,0}(z)=\kappa_e u_*z.
 \end{aligned}}
 \label{ba:outer-profile}
\end{equation}
The velocity can be written as $u_*\ln(z/z_{0,w})/\kappa_e$, where
\begin{equation}
 \boxed{z_{0,w}=H\exp[-\kappa_e u_*J_0(H)].}
 \label{ba:aerodynamic-intercept}
\end{equation}
Thus fixed $h$, $H$, and viscosity parameters determine an aerodynamic
intercept. A separately prescribed $z_0$ must satisfy
\begin{equation}
 \ln(H/z_0)=\kappa_e u_*J_0(H).
 \label{ba:prescribed-z0}
\end{equation}
It is not permissible to prescribe an incompatible $z_0$ and retain the
same AWF momentum anchor while claiming to preserve the same complete
reference. The roughness/intercept calibration needed for a prescribed
$z_0$ is distinct from the turbulence-budget calibration.

\subsection{Budget-determined dissipation and admissibility}
\label{ba:cutoff}

Choose the continuous positive inner sink
\begin{equation}
 \boxed{E_w(z)=\frac{\mathcal A}{\max(z,z_c)},\qquad C_k^w(z)=0,}
 \label{ba:dissipation}
\end{equation}
where $C_k^w$ would denote an additional prescribed redistribution source.
The tail must begin no later than the matching height, so $0<z_c\leq H$.
Then
\begin{equation}
 \int_0^HE_w\,\mathrm dz
 =\mathcal A\left(1+\ln\frac H{z_c}\right).
 \label{ba:dissipation-integral}
\end{equation}
The outer reference has zero $k$ gradient. Conservative matching to that
reference and zero modeled ground $k$ flux therefore require
\begin{equation}
 \int_0^HE_w\,\mathrm dz=\mathcal P_0.
 \label{ba:budget-condition}
\end{equation}
Solving this scalar condition gives
\begin{equation}
 \boxed{z_c=H\exp(1-\Lambda),\qquad
 \Lambda=\frac{\mathcal P_0}{\mathcal A}
 =\frac{t_0^2[J_0(H)-\nu J_2(H)]}{\mathcal A}.}
 \label{ba:cutoff-formula}
\end{equation}
The left-hand side of \eqref{ba:dissipation-integral} decreases strictly
with $z_c$, diverges as $z_c\downarrow0$, and equals $\mathcal A$ at
$z_c=H$. Consequently,
\begin{equation}
 \boxed{\text{A unique }0<z_c\leq H\text{ exists if and only if }\Lambda\geq1.}
 \label{ba:cutoff-gate}
\end{equation}
If a separate comparison requires the $1/z$ tail from an interior height
$z_m<H$, the corresponding stronger gate is
\begin{equation}
 z_c\leq z_m\quad\Longleftrightarrow\quad
 \Lambda\geq1+\ln(H/z_m).
 \label{ba:interior-gate}
\end{equation}
Failure of the applicable gate excludes the reference from this monotone
cutoff family. Clipping $z_c$ to $H$ or $z_m$ would destroy the budget.

The original continuity-based molecular cutoff, with the calibrated
coefficient inserted, would be
\begin{equation}
 z_{\epsilon,A}=\frac{2c_\ell^w\nu}{\sqrt K},\qquad
 E_A(z)=\frac{\mathcal A}{\max(z,z_{\epsilon,A})}.
 \label{ba:original-cutoff}
\end{equation}
The replacement \eqref{ba:cutoff-formula} is a change of the dissipation
model, not a new evaluation of that original molecular cutoff. Its
plateau equals $\mathcal A/z_c$. Imposing the additional plateau formula
$2\nu K/z_c^2$ would force $z_c=z_{\epsilon,A}$ again, leaving no
independent budget adjustment.

\subsection{The nonconstant inner reference and all test-function moments}
\label{ba:inner-k}

Retain inner transport with actual production and the prescribed sink:
\begin{equation}
 \partial_t k_w+\boldsymbol U_w\cdot\nabla k_w
 =\nabla\cdot(D_{k,w}\nabla k_w)+P_A-E_w.
 \label{ba:k-transport}
\end{equation}
For the flat, steady, parallel reference the following derivation works
for any prescribed positive normal diffusivity $D_z(z)$. The selected
model sets $D_z=D_{k,w}=D_A$. The boundary-value problem is
\begin{equation}
 (D_zk_{w,0}')'=E_w-p_{A,0},\qquad
 D_z(0)k_{w,0}'(0)=0,\qquad k_{w,0}(H)=K.
 \label{ba:k-bvp}
\end{equation}
Define the cumulative source defect
\begin{equation}
 Q_0(z)=\int_0^z[E_w(s)-p_{A,0}(s)]\,\mathrm ds.
 \label{ba:Q}
\end{equation}
Two integrations give the unique solution for that prescribed diffusivity:
\begin{equation}
 \boxed{D_z(z)k_{w,0}'(z)=Q_0(z),\qquad
 k_{w,0}(z)=K-\int_z^H\frac{Q_0(s)}{D_z(s)}\,\mathrm ds.}
 \label{ba:k-profile}
\end{equation}
$Q_0$ is the gradient flux inside the diffusion operator; the physical
upward diffusive flux is $-Q_0$. The budget implies
\begin{equation}
 Q_0(0)=Q_0(H)=0,\qquad k_{w,0}'(H)=0.
 \label{ba:endpoint-fluxes}
\end{equation}
Conversely, zero gradient flux at both endpoints implies the integrated
budget when \eqref{ba:k-bvp} is integrated. Allowing a nonconstant inner
$k$ supplies local diffusion but does not eliminate this compatibility
condition.

No additional value is prescribed for $k_w(0)$: it is determined by
\eqref{ba:k-profile}. This is a modeled rough-wall layer, with no slip for
the mean velocity, rather than a microscopic viscous-sublayer resolution
with an independently imposed turbulent condition $k(0)=0$. The formula
defines a reference and an initialization. During a general flow, the
field evolves under \eqref{ba:k-transport} and is not reset to that profile.

For every sufficiently regular test function $\phi$, integration by parts
gives
\begin{equation}
 \boxed{-\int_0^H D_zk_{w,0}'\phi'\,\mathrm dz
       +\int_0^H(p_{A,0}-E_w)\phi\,\mathrm dz=0.}
 \label{ba:weak-equilibrium}
\end{equation}
The result balances every continuum test moment, including nonconstant
ones. A constant-inner-$k$ construction without redistribution would
lose the first term and could balance only those moments for which the
source integral happened to vanish. Projection, quadrature, and numerical
differentiation can still leave a discrete residual; the continuum
identity does not by itself establish exact spectral/$hp$ or finite-volume
preservation.

\subsection{Closed formulas and positivity on the rough branch}
\label{ba:rough}

Assume $z_v\leq0$, and define
\begin{equation}
 d=\nu-az_v\geq\nu,
 \qquad D_A(z)=d+az.
 \label{ba:rough-D}
\end{equation}
Then
\begin{equation}
 \begin{aligned}
 J_0(z)&=\frac1a\ln\frac{d+az}{d},&
 J_2(z)&=\frac1a\left(\frac1d-\frac1{d+az}\right),\\
 p_{A,0}(z)&=t_0^2\frac{d+az-\nu}{(d+az)^2}.&&
 \end{aligned}
 \label{ba:rough-integrals}
\end{equation}
The cumulative production is
\begin{equation}
 \int_0^zp_{A,0}(s)\,\mathrm ds
 =\mathcal A\left[
 \ln\frac{d+az}{d}
 -\nu\left(\frac1d-\frac1{d+az}\right)\right].
 \label{ba:rough-production}
\end{equation}
With $\zeta_H=aH/d$ and $\beta=\nu/d$,
\begin{equation}
 \Lambda=\ln(1+\zeta_H)-\frac{\beta\zeta_H}{1+\zeta_H}.
 \label{ba:rough-Lambda}
\end{equation}
Define the dimensionless cumulative sink
\begin{equation}
 R_c(z)=\begin{cases}
 z/z_c,&0\leq z\leq z_c,\\
 1+\ln(z/z_c),&z_c\leq z\leq H.
 \end{cases}
 \label{ba:Rc}
\end{equation}
Then the flux has the elementary representation
\begin{equation}
 \boxed{Q_0(z)=\mathcal A\left[
 R_c(z)-\ln\frac{d+az}{d}
 +\nu\left(\frac1d-\frac1{d+az}\right)\right].}
 \label{ba:rough-Q}
\end{equation}
Together with \eqref{ba:k-profile}, this is already an exact analytical
representation of the inner reference.

\paragraph{A closed primitive for $k_{w,0}$.}
For $\sigma_k=1$, put $\zeta=az/d$, $\zeta_c=az_c/d$, and define
\begin{align}
 \mathcal F_< (\zeta)
 &=\frac{\zeta-\ln(1+\zeta)}{\zeta_c}
   -\frac12\ln^2(1+\zeta)
   +\beta\left[\ln(1+\zeta)+\frac1{1+\zeta}-1\right],
 \label{ba:primitive-lower}\\
 \mathcal G(\zeta)
 &=\left[1+\ln\frac{\zeta}{\zeta_c}+\beta\right]\ln(1+\zeta)
   +\operatorname{Li}_2(-\zeta)
   -\frac12\ln^2(1+\zeta)+\frac{\beta}{1+\zeta},\qquad\zeta>0.
 \label{ba:primitive-upper}
\end{align}
The real dilogarithm is
$\operatorname{Li}_2(-\zeta)=-\int_0^\zeta\ln(1+s)\,\mathrm ds/s$.
Join the primitives continuously:
\begin{equation}
 \mathcal F(\zeta)=
 \begin{cases}
 \mathcal F_< (\zeta),&0\leq\zeta\leq\zeta_c,\\
 \mathcal F_< (\zeta_c)+\mathcal G(\zeta)-\mathcal G(\zeta_c),
 &\zeta_c\leq\zeta\leq\zeta_H.
 \end{cases}
 \label{ba:primitive-joined}
\end{equation}
Direct differentiation gives
$\mathcal F'(\zeta)=Q_0(d\zeta/a)/[\mathcal A(1+\zeta)]$, hence
\begin{equation}
 \boxed{k_{w,0}(z)=K+\frac{\mathcal A}{a}
 \left[\mathcal F\!\left(\frac{az}{d}\right)
       -\mathcal F(\zeta_H)\right].}
 \label{ba:k-closed}
\end{equation}
The integral solution \eqref{ba:k-profile} remains valid on other
viscosity segments and for other prescribed positive normal diffusivities.

\paragraph{The sign of the local source defect.}
For $z\geq z_c>0$ on this rough branch,
\begin{equation}
 \boxed{E_w(z)-p_{A,0}(z)
 =\frac{\mathcal A[(d+\nu)az+d^2]}{z(d+az)^2}>0.}
 \label{ba:tail-defect}
\end{equation}
This strict inequality applies to the $1/z$ tail. The zero total budget
requires a compensating negative defect somewhere on the plateau.

\paragraph{Proposition: a positive rough-wall reference.}
Suppose $z_v\leq0$, $d\geq2\nu$, and $\Lambda\geq1$. Then
$E_w-p_{A,0}$ has exactly one zero in $(0,z_c)$ and changes there from
negative to positive. For every prescribed positive $D_z$,
\begin{equation}
 \boxed{Q_0(z)<0\ (0<z<H),\qquad
 k_{w,0}'(z)<0\ (0<z<H),\qquad
 k_{w,0}(z)>K>0\ (0\leq z<H).}
 \label{ba:positivity}
\end{equation}
\textit{Proof.}
The reference production satisfies
\[
 p_{A,0}'(z)
 =\frac{a t_0^2[2\nu-(d+az)]}{(d+az)^3}\leq0,
\]
with strict decrease for $z>0$. On the plateau,
$\mathcal A/z_c-p_{A,0}$ is strictly increasing and is positive at $z_c$
by \eqref{ba:tail-defect}. Its value at zero must be negative, since
otherwise the full-layer integral could not vanish. It therefore crosses
zero exactly once. The defect stays positive after the crossing, so its
cumulative integral decreases first and then increases to $Q_0(H)=0$,
remaining negative in the interior. Formula \eqref{ba:k-profile} proves
the claims.\hfill$\square$

The condition $d\geq2\nu$ is sufficient, not necessary. For other
intercepts, including a smooth-wall viscosity segment, the cutoff remains
defined when its gate holds, but positivity must be checked from the
resulting profile.

\paragraph{Covariance realizability of this reference.}
Use the full-strain auxiliary stress completion
\begin{equation}
 R_w=\frac23 k_w I-\nu_{t,A}\mathsf M_w,
 \qquad \mathsf M_w=\nabla\boldsymbol U_w+
                         (\nabla\boldsymbol U_w)^T.
 \label{ba:inner-stress-definition}
\end{equation}
For the flat reference, put
$\tau_t(z)=\nu_{t,A}U_{w,0}'=t_0\nu_{t,A}/D_A$. Its covariance
eigenvalues are
\begin{equation}
 \lambda(R_w)=
 \left\{\frac23k_{w,0},\quad
        \frac23k_{w,0}+|\tau_t|,\quad
        \frac23k_{w,0}-|\tau_t|\right\}.
 \label{ba:inner-covariance}
\end{equation}
Under the proposition's assumptions, $k_{w,0}\geq K=q_et_0$ and
$0\leq\tau_t<t_0$, since $\nu>0$. Thus
\begin{equation}
 \boxed{\lambda_{\min}(R_w)
 >\left(\frac{2q_e}{3}-1\right)t_0>0.}
 \label{ba:reference-realizability}
\end{equation}
On the rough branch $\tau_t$ increases while $k_{w,0}$ decreases, so
the minimum eigenvalue occurs at $H$. This proves covariance
realizability of the constructed flat reference, rather than inferring it
from $k>0$ alone. It is not a realizability theorem for the auxiliary
stress in arbitrary evolving strain fields, nor a claim about incremental
constitutive stability.

\paragraph{Aerodynamic intercept and cutoff.}
The rough-branch anchor gives
\begin{equation}
 \boxed{z_{0,w}=\frac{dH}{aH+d},\qquad
 z_c=\mathrm e\,z_{0,w}
 \exp\!\left[\nu\left(\frac1d-\frac1{aH+d}\right)\right].}
 \label{ba:roughness-cutoff}
\end{equation}
The aerodynamic intercept retains a finite-height dependence even with
the calibrated slope; it tends to $d/a$ as $aH/d\to\infty$.
For a prescribed $0<z_0<H$, the rough-branch momentum calibration requires
\begin{equation}
 d=\frac{aHz_0}{H-z_0},\qquad z_v=\frac{\nu-d}{a}.
 \label{ba:inverse-intercept}
\end{equation}
An empirical roughness value producing that intercept must be admissible
under \eqref{ba:roughness-correlation}; the cutoff and positivity gates
still require separate checks.

\subsection{Conservative matching, pressure balance, and mechanical energy}
\label{ba:matching}

The inner and outer reference fields satisfy
\begin{equation}
 \begin{gathered}
 U_{w,0}(0)=0,\qquad U_{w,0}(H)=U_{o,0}(H),\\
 D_A(H)U_{w,0}'(H)
 =-R_{xz,o,0}(H)
 =\nu_{t,o,0}(H)U_{o,0}'(H)=t_0.
 \end{gathered}
 \label{ba:momentum-match}
\end{equation}
Equal slopes of the viscosity models do not mean equal viscosities at
finite $H$: $D_A(H)=aH+d$, whereas the high-Re outer reference has
$\nu_{t,o,0}(H)=aH$. The velocity gradients need not be continuous;
velocity and total tangential traction are continuous.

The turbulence conditions are
\begin{equation}
 \boxed{\begin{aligned}
 k_{w,0}(H)&=k_{o,0}(H)=K,\\
 D_{k,w}(H)k_{w,0}'(H)&=D_{k,o}(H)k_{o,0}'(H)=0,\\
 E_w(H)&=\epsilon_{o,0}(H)=\frac{\mathcal A}{H}.
 \end{aligned}}
 \label{ba:turbulence-match}
\end{equation}
Here $D_{k,o}=C_dk_o^2/(\sigma_k\epsilon_o)$ belongs to the declared
outer equations. Inner dissipation is dependent, so the outer
dissipation equation receives this single trace and determines its flux.
No additional dissipation-flux condition is imposed. At the limiting
cutoff $z_c=H$, the one-sided dissipation derivatives differ, although the
required trace still matches.

The outer WJ reference normal stress is
\begin{equation}
 R_{zz,o,0}=\left(\frac23-\frac1{N_e}\right)K
 =\frac{34}{81}K,
 \label{ba:outer-normal-stress}
\end{equation}
whereas \eqref{ba:inner-stress-definition} gives
$R_{zz,w,0}=2k_{w,0}/3$. Inner normal momentum therefore requires
\begin{equation}
 \frac{\mathrm d}{\mathrm dz}
 \left(\frac{p_w}{\rho}+gz+\frac23k_{w,0}\right)=0.
 \label{ba:inner-pressure}
\end{equation}
Mechanical pressure is generally not hydrostatic inside this layer.
Continuity of total normal traction at the flat interface gives
\begin{equation}
 \boxed{\frac{p_w(H)-p_o(H)}{\rho}
 =-\frac{K}{N_e}=-\frac{20}{81}K.}
 \label{ba:pressure-offset}
\end{equation}
The offset depends on the matching value $K$; the interior pressure
distribution also depends on the nonconstant $k_{w,0}$. If the solver
stores a reduced pressure, mechanical pressure must be recovered with the
correct local isotropic stress and gravitational potential before forming
the sampled AWF pressure forcing. The pressure representation must allow
the mechanical-pressure offset, or impose continuity of an appropriately
defined total-normal-stress pressure instead. Enforcing a common mechanical
pressure at this interface would contradict \eqref{ba:pressure-offset}.

The flat reference obeys the integrated mechanical-energy identity
\begin{equation}
 \boxed{t_0U_{H,0}
 =\nu\int_0^H(U_{w,0}')^2\,\mathrm dz
  +\int_0^Hp_{A,0}\,\mathrm dz
 =\nu t_0^2J_2(H)+\int_0^HE_w\,\mathrm dz.}
 \label{ba:energy}
\end{equation}
The nonzero internal $k$ flux redistributes energy within the modeled
layer but contributes no net flux at the two reference endpoints.

\subsection{Comparison with an EP-AWF at unit equilibrium weight}
\label{ba:ep-comparison}

For a precise comparison, use identical $K,h,H,\nu$, calibrated viscosity
coefficients, and the same admissible budget cutoff in both constructions.
An EP-AWF at unit equilibrium weight can retain constant inner $k=K$ by
using
\begin{equation}
 E_w^{EP}=E_w,\qquad
 C_{k,EP}^w=E_w-p_{A,0},\qquad
 S_{k,w}^{EP}=P_A-E_w^{EP}+C_{k,EP}^w=P_A-p_{A,0}.
 \label{ba:ep-source}
\end{equation}
At the flat reference, this local source vanishes pointwise. The
budget-cutoff variant instead uses
\begin{equation}
 E_w=\frac{\mathcal A}{\max(z,z_c)},\qquad
 C_k^w=0,\qquad S_{k,w}=P_A-E_w,
 \label{ba:budget-source}
\end{equation}
whose nonzero reference source is balanced by
$\partial_z(D_zk_{w,0}')$.

For the same instantaneous $P_A$ and sample $K$, the layer integrals agree:
\begin{equation}
 \int_0^HS_{k,w}\,\mathrm dz
 =\int_0^HS_{k,w}^{EP}\,\mathrm dz
 =\int_0^H(P_A-p_{A,0})\,\mathrm dz.
 \label{ba:equal-integrated-sources}
\end{equation}
Their local balances and transported fields differ. At the reference,
the EP redistribution can be written as $\partial_zQ_0$ and carries
physical flux $-Q_0$; ordinary diffusion of the variable profile carries
the same reference flux in the budget-cutoff variant. This identity
does not make the two evolving models dynamically equivalent.

If an EP transition is defined by
\begin{equation}
 E_w^{EP}(\chi)=(1-\chi)E_A+\chi E_w,\qquad
 C_{k,EP}^w(\chi)=\chi(E_w-p_{A,0}),
 \label{ba:ep-transition}
\end{equation}
then setting $\chi=0$ returns the molecular-cutoff law $E_A$, not the
budget-cutoff variant. The latter retains $E_w$ while removing the
redistribution. It is therefore a distinct model selection, with its
own reference profile and positivity assessment.

\subsection{A conservative extension to curved terrain}
\label{ba:curved}

The following extension retains the calibrated coefficient laws and the
source policy $P_A-E_w$, but evaluates the reference cutoff using physical
volume measure. It is a specified conservative extension; the flat
equilibrium result does not assert that flow over a hill remains logarithmic.

\paragraph{Geometry and admissibility.}
For a stationary smooth wall $\boldsymbol X(\boldsymbol\xi)$, let the unit
normal $\boldsymbol n$ point into the fluid. Use
\begin{equation}
 \begin{gathered}
 \boldsymbol x(\boldsymbol\xi,r)
 =\boldsymbol X(\boldsymbol\xi)+r\boldsymbol n(\boldsymbol\xi),\qquad
 0\leq r\leq H(\boldsymbol\xi),\\
 P_T=I-\boldsymbol n\otimes\boldsymbol n,\qquad
 \mathsf B=\nabla_\Gamma\boldsymbol n,\qquad
 \mathsf A=I_T+r\mathsf B,\qquad J=\det_T\mathsf A,\\
 \mathrm dV=J\,\mathrm dS_w\,\mathrm dr,\qquad
 \nabla f=\mathsf A^{-T}\nabla_\Gamma f+\boldsymbol n\partial_rf.
 \end{gathered}
 \label{ba:curved-map}
\end{equation}
The height field is fixed in time. Surface derivatives hold $r$ fixed,
and vectors retain Cartesian components. In this sign convention an
outward-oriented sphere has positive principal curvatures. All stretches
$1+r\kappa_i$ must remain positive; the normal map must also be globally
one-to-one, since $J>0$ alone does not exclude intersecting normal rays.
Numerical geometry and interface measures must be compatible. Require
$0\leq h<H$, $K>0$, and sample $K$ from the common trace at $r=H$.

\paragraph{Reference cutoff in physical volume measure.}
Along each normal ray, define
\begin{equation}
 \boxed{\begin{aligned}
 E_w(r)&=\frac{\mathcal A(K)}{\max(r,r_c)},\\
 \int_0^HJ(r)E_w(r)\,\mathrm dr
 &=\mathcal P_J:=\int_0^HJ(r)p_{A,0}(r;K)\,\mathrm dr,\\
 p_{A,0}(r;K)&=t_0(K)^2
                 \frac{\nu_{t,A}(r;K)}{D_A(r;K)^2}.
 \end{aligned}}
 \label{ba:curved-budget}
\end{equation}
The reference kernel remains the counterfactual flat, parallel,
zero-forcing production at the current $K$. It is not the production
of the current curved flow.
For $\mathcal B(c)=\int_0^HJ\mathcal A/\max(r,c)\,\mathrm dr$,
\begin{equation}
 \mathcal B'(c)=-\frac{\mathcal A}{c^2}\int_0^cJ(r)\,\mathrm dr<0,
 \qquad \lim_{c\downarrow0}\mathcal B(c)=+\infty.
 \label{ba:curved-monotonicity}
\end{equation}
A unique admissible cutoff exists precisely when
\begin{equation}
 \boxed{\mathcal P_J\geq
 \frac{\mathcal A}{H}\int_0^HJ(r)\,\mathrm dr.}
 \label{ba:curved-gate}
\end{equation}
At a fixed footpoint, $J=(1+\kappa_1r)(1+\kappa_2r)$, and direct
integration gives
\begin{equation}
 \frac{\mathcal B(c)}{\mathcal A}
 =1+\ln\frac Hc
 +(\kappa_1+\kappa_2)\left(H-\frac c2\right)
 +\kappa_1\kappa_2\left(\frac{H^2}{2}-\frac{c^2}{6}\right).
 \label{ba:curved-budget-explicit}
\end{equation}
Thus one monotone scalar root solve replaces the flat exponential formula.

\paragraph{Actual flow equations.}
Write $\boldsymbol U=\boldsymbol u_t+w\boldsymbol n$ and
$\boldsymbol v=\mathsf A^{-1}\boldsymbol u_t$. Continuity and momentum
are
\begin{align}
 &\partial_r(Jw)+\operatorname{div}_\Gamma(J\boldsymbol v)=0,
 \label{ba:curved-continuity}\\
 &J\partial_t\boldsymbol U
 +\partial_r\{J(w\boldsymbol U-\mathsf T_w\boldsymbol n)\}
 +\operatorname{Div}_\Gamma
    \{J(\boldsymbol U\otimes\boldsymbol v-\mathsf T_w\mathsf A^{-T})\}
 =J\boldsymbol f.
 \label{ba:curved-momentum}
\end{align}
Here $\operatorname{Div}_\Gamma$ acts on Cartesian rows, and the
full-strain completion is
\begin{equation}
 \begin{gathered}
 R_w=\tfrac23k_wI-\nu_{t,A}\mathsf M_w,\qquad
 \mathsf T_w=-\left(\frac{p_w}{\rho}+\tfrac23k_w\right)I+D_A\mathsf M_w,\\
 P_A=-R_w:\nabla\boldsymbol U_w
     =\tfrac12\nu_{t,A}\mathsf M_w:\mathsf M_w
       \quad\text{when }\nabla\cdot\boldsymbol U_w=0.
 \end{gathered}
 \label{ba:curved-stress}
\end{equation}
The isotropic term uses local $k_w$, not sampled $K$. If the discrete
field is only weakly incompressible, the full contraction additionally
contains $-(2k_w/3)\nabla\cdot\boldsymbol U_w$; retaining that contraction
with the same momentum stress preserves the chosen discrete stress-work
accounting.

One definite representation of the original roughness forcing suppression
is
\begin{equation}
 \begin{aligned}
 \boldsymbol f&=\boldsymbol b+
           (1-\mathbf1_{r>h})\boldsymbol F_A,\\
 \boldsymbol F_A&=\mathcal E_H^+P_T
 \left[\operatorname{Div}_r
           (\boldsymbol u_{o,t}\otimes\boldsymbol u_{o,t})
        +\nabla_r(p_o/\rho)-\boldsymbol b\right],\\
 \nabla_r&=\mathsf A^{-T}\nabla_\Gamma,\qquad
 \operatorname{Div}_r\mathsf Q
 =J^{-1}\operatorname{Div}_\Gamma(J\mathsf Q\mathsf A^{-T})
 \quad(\mathsf Q\boldsymbol n=0).
 \end{aligned}
 \label{ba:rough-force}
\end{equation}
$\boldsymbol b$ contains physical body acceleration. The one-sided
sampler $\mathcal E_H^+$ evaluates outer quantities at fixed $r$ before
restriction to the graph $r=H$ and extends the sample constantly along
the ray. Hydrostatic pressure and gravity must cancel in this sampler.
The additional roughness force and its power belong in the budgets.
Under the original flat, quasi-steady sampled-forcing approximation,
\eqref{ba:curved-momentum} reduces to \eqref{ba:momentum-ode}.

The actual inner kinetic-energy equation is
\begin{equation}
 \boxed{\begin{aligned}
 J\partial_t k_w
 &+\partial_r\{J(wk_w-D_{k,w}\partial_rk_w)\}\\
 &+\operatorname{div}_\Gamma\left\{
 J[\boldsymbol v k_w-D_{k,w}\mathsf A^{-1}\mathsf A^{-T}
                                 \nabla_\Gamma k_w]\right\}
 =J(P_A-E_w).
 \end{aligned}}
 \label{ba:curved-k}
\end{equation}
The base-normal tangential traction contains curvature and normal motion:
\begin{equation}
 \boldsymbol t_A=D_A\left[
 \partial_r\boldsymbol u_t+\mathsf A^{-T}\nabla_\Gamma w
 -\mathsf B\mathsf A^{-1}\boldsymbol u_t\right].
 \label{ba:curved-traction}
\end{equation}
Normal momentum retains curvature acceleration and normal-stress gradients
through \eqref{ba:curved-momentum}. Weighting independent flat columns
by $J$ does not supply these terms; general terrain flow requires the
coupled patch equations.

The reference budget preserves the integrated actual departure from the
reference:
\begin{equation}
 \int_0^HJ(P_A-E_w)\,\mathrm dr
 =\int_0^HJ(P_A-p_{A,0})\,\mathrm dr.
 \label{ba:actual-budget}
\end{equation}
This quantity can be nonzero and is balanced by storage, tangential
transport, and interface flux. Choosing the cutoff from actual $P_A$
instead of reference $p_{A,0}$ would impose a different energy budget.

An auxiliary frozen normal-column construction remains available:
\begin{equation}
 \begin{aligned}
 Q_J(r)&=\int_0^rJ(s)[E_w(s)-p_{A,0}(s)]\,\mathrm ds,\\
 k_{\mathrm{col}}(r)&=K-\int_r^H
               \frac{Q_J(s)}{J(s)D_{k,w}(s)}\,\mathrm ds.
 \end{aligned}
 \label{ba:curved-column}
\end{equation}
It balances reference production and normal diffusion only. Even when
its endpoint gradient flux is zero, it is not an exact terrain-flow
equilibrium, because actual production, advection, and tangential
diffusion generally differ. A zero base-normal gradient also need not
give zero diffusive flux through a sloping matching surface.

\paragraph{Conservative interface conditions.}
For spatially varying $H$, put
\begin{equation}
 \boldsymbol q_H=\mathsf A_H^{-T}\nabla_\Gamma H,\qquad
 \beta_H=\sqrt{1+|\boldsymbol q_H|^2},\qquad
 \boldsymbol N_H=\frac{\boldsymbol n-\boldsymbol q_H}{\beta_H},\qquad
 \mathrm dS_H=J_H\beta_H\,\mathrm dS_w.
 \label{ba:interface-geometry}
\end{equation}
Here $\boldsymbol N_H$ points into the outer region. The mass flux per
unit ground area is
\begin{equation}
 m_H=J_H(w_H-\boldsymbol v_H\cdot\nabla_\Gamma H)
 =-\operatorname{div}_\Gamma\int_0^HJ\boldsymbol v\,\mathrm dr.
 \label{ba:interface-mass}
\end{equation}
Use one common velocity trace, mass flux, and total traction:
\begin{equation}
 \boldsymbol U_w=\boldsymbol U_o=\widehat{\boldsymbol U},\qquad
 \mathsf T_w\boldsymbol N_H=\mathsf T_o\boldsymbol N_H,
 \qquad \mathsf T_o=-\frac{p_o}{\rho}I-R_o,
 \label{ba:interface-momentum}
\end{equation}
where $R_o$ is the full outer WJ covariance. Inner and outer weak
equations use opposite orientations of the same advective and stress
fluxes and the same interface measure. The normal part gives
\begin{equation}
 \frac{p_w-p_o}{\rho}
 =D_A\boldsymbol N_H\cdot\mathsf M_w\boldsymbol N_H
   -\frac23k_w+\boldsymbol N_H\cdot R_o\boldsymbol N_H,
 \label{ba:curved-pressure}
\end{equation}
which replaces the flat offset \eqref{ba:pressure-offset}.

Require a common $k$ trace and a common advective--diffusive flux. For
$\boldsymbol f_k=\boldsymbol U k-D_k\nabla k$, its value per unit
ground area is
\begin{equation}
 J_H\beta_H\boldsymbol f_k\cdot\boldsymbol N_H
 =m_H\widehat k-J_HD_k
  \left[\partial_rk-
        \boldsymbol q_H\cdot\mathsf A_H^{-T}\nabla_\Gamma k\right].
 \label{ba:interface-k}
\end{equation}
Ground modeled $k$ flux is zero. For the endpoint sample,
$\nabla_\Gamma K=(\nabla_\Gamma k)|_H+
(\partial_rk)|_H\nabla_\Gamma H$; surface derivatives at fixed $r$ must
therefore be distinguished from derivatives of the matching trace.
The outer dissipation equation receives only
\begin{equation}
 \boxed{\epsilon_{o,H}=E_w(H)=\frac{\mathcal A(K)}{H}.}
 \label{ba:epsilon-trace}
\end{equation}
There is no additional inner dissipation equation or prescribed
dissipation flux. The conservative coupling claims concern mass,
momentum, and the retained kinetic energy.

\subsection{A first-zero obstruction and an optional positivity amendment}
\label{ba:dynamic-positivity}

\paragraph{Why the default sampled sink can produce negative $k_w$.}
The positive reference in \eqref{ba:positivity} does not imply that
the default evolving equation \eqref{ba:k-transport} preserves
nonnegativity. Its sink is controlled by the endpoint sample $K$ and
does not vanish when the local transported value vanishes.
For an explicit boundary-of-the-positive-cone test, fix any admissible
$K>0$ and take zero actual mean velocity and forcing, so $P_A=0$.
Let
\begin{equation}
 \psi(s)=
 \begin{cases}
 0,&s\leq0,\\
 10s^3-15s^4+6s^5,&0<s<1,\\
 1,&s\geq1,
 \end{cases}
 \qquad k_w(z,0)=K\psi(4z/H-2).
 \label{ba:first-zero-data}
\end{equation}
This $C^2$ nonnegative field has zero ground flux, is identically zero
below $H/2$, and equals $K$ near $H$. In particular, endpoint sampling
still gives the prescribed positive $K$. The resting mean flow satisfies
no slip, and its normal momentum can be balanced at this instant by
$p_w/\rho+gz+2k_w/3=\mathrm{const}$.
At $z_*=H/4$ both advection
and diffusion of $k_w$ vanish, while the counterfactual budget still
gives $E_w(z_*)>0$. Therefore
\begin{equation}
 \left.\partial_t k_w\right|_{z_*,\,t=0}
 =-E_w(z_*)<0.
 \label{ba:first-zero-failure}
\end{equation}
Any regular evolution from this state immediately becomes negative
there. Equivalently, the same local configuration at a first zero has
a destruction term pointing out of the nonnegative cone. The source
already fails the necessary quasi-positivity test before discretization.
The reference-production budget cannot repair this failure because
reference production is not actual production.

\paragraph{A conditional modification with a bounded destruction rate.}
The following is an optional modified variant. It introduces a
redistribution and, under severe depletion, reduces the dissipation
budget. It is not part of the default $C_k^w=0$ formulation.
Use $J=1$ on a flat layer or the positive geometric Jacobian in a curved
layer, and define
\begin{equation}
 \mathcal B_w=\int_0^H J E_w\,\mathrm dr>0,\qquad
 0<\theta<1,\qquad 0<\delta<1.
 \label{ba:positive-budget}
\end{equation}
For nonnegative local $k_w$ and positive sampled $K$, set
\begin{equation}
 \omega_k=\min\{1,k_w/(\theta K)\},\qquad
 b_k=\frac{\int_0^H J E_w\omega_k\,\mathrm dr}{\mathcal B_w},
 \qquad m_k=\max(b_k,\delta),\qquad \lambda_k=\frac{b_k}{m_k}.
 \label{ba:positive-weights}
\end{equation}
The weighted mean $b_k$ is unrelated to the roughness intercept $b$.
The destruction, actual dependent dissipation, and redistribution are
\begin{equation}
 \boxed{D_+=\frac{E_w\omega_k}{m_k},\qquad
 E_+=\lambda_kE_w,\qquad C_+=E_+-D_+.}
 \label{ba:positive-profiles}
\end{equation}
$D_+$ has dissipation-rate units and is not a diffusivity.
Replace the scalar source by
$P_A-E_++C_+=P_A-D_+$, keeping the stated positive diffusivity.
For $k_w>0$,
\begin{equation}
 0\leq\frac{D_+}{k_w}\leq\frac{E_w}{\delta\theta K},
 \qquad D_+|_{k_w=0}=0.
 \label{ba:positive-rate}
\end{equation}
Thus the reaction source at a first zero is $P_A\geq0$.
With regular solutions, bounded admissible coefficients, positive
diffusivity, and compatible nonnegative boundary data, this removes
the source obstruction in \eqref{ba:first-zero-failure}. It does not
prove positivity of the coupled outer equations or of a high-order
space--time discretization. For a weakly incompressible velocity,
the isotropic production contribution also vanishes at $k_w=0$;
the stress-work convention must remain consistent away from that state.

The exact integral identities are
\begin{equation}
 \boxed{\int_0^H JD_+\,\mathrm dr
 =\int_0^H JE_+\,\mathrm dr=\lambda_k\mathcal B_w,\qquad
 \int_0^H JC_+\,\mathrm dr=0.}
 \label{ba:positive-integrals}
\end{equation}
If $\omega_k=1$ everywhere, then $b_k=m_k=\lambda_k=1$:
$D_+=E_+=E_w$ and $C_+=0$ pointwise. In particular, the rough
reference obeys $k_{w,0}\geq K>\theta K$, so the amendment is exactly
inactive, including in an open neighborhood of that reference.
Its nonconstant reference profile, original cutoff budget, pressure
balance, and all reference test moments are preserved.

When the limiter is active and $b_k\geq\delta$, the original integrated
budget and dependent dissipation are retained, but destruction is
redistributed away from depleted regions. Healthy regions can receive
larger destruction through the factor $1/b_k$; this is not a pointwise
reduction everywhere. When $b_k<\delta$, the cap keeps the local
destruction rate bounded by reducing the integrated loss to
$(b_k/\delta)\mathcal B_w$. Such a concession is necessary if both
a bounded rate and severe depletion are admitted: a bound
$D_+\leq a_{\max}k_w$ implies
$\int JD_+\leq a_{\max}\int Jk_w$, which cannot sustain a fixed
positive loss as available layer energy tends to zero.

The redistribution has the conservative representation
\begin{equation}
 \boldsymbol j_+(r)=\frac{\boldsymbol n}{J(r)}
       \int_0^rJ(s)C_+(s)\,\mathrm ds,\qquad
 \nabla\cdot\boldsymbol j_+=C_+,\qquad
 \boldsymbol j_+(0)=\boldsymbol j_+(H)=0.
 \label{ba:positive-flux}
\end{equation}
One may use physical flux
$\boldsymbol U k_w-D_{k,w}\nabla k_w-\boldsymbol j_+$
with source $P_A-E_+$, or ordinary diffusion with source $P_A-D_+$.
These are alternative representations and must not both be assembled.
The interface $k$ flux receives no additional term because
$\boldsymbol j_+(H)=0$.

The outer dissipation trace must follow the actual dependent field:
\begin{equation}
 \boxed{\epsilon_{o,H}=E_+(H)
 =\lambda_k\frac{\mathcal A(K)}H.}
 \label{ba:positive-trace}
\end{equation}
Retaining the unscaled trace when $\lambda_k<1$ would be inconsistent
with the declared inner dissipation. For continuous endpoint-sampled
$k_w$ and $K>0$, each regular state has $b_k>0$, but a uniform lower
bound is not established as its positive support contracts. The outer
closure's $k/\epsilon$ ratios consequently need a separately qualified
continuation if that degeneracy is approached. No undisclosed
dissipation floor is included here. Metric normalization, cutoff,
sample $K$, and the trace must be recomputed consistently after
state-changing projection or limiting; positive source quadrature
including the endpoint prevents an exact zero normalization at fixed
resolution, without establishing a uniform bound under refinement.

\subsection{Quantitative reference and implementation implications}
\label{ba:example}

Use the flat reference inputs
\begin{equation}
 u_*=0.5\,\mathrm{m\,s^{-1}},\qquad
 \nu=1.5\times10^{-5}\,\mathrm{m^2\,s^{-1}},\qquad
 h=0.1\,\mathrm m,\qquad H=1\,\mathrm m.
 \label{ba:example-inputs}
\end{equation}
The unchanged empirical intercept correlation gives
$h^*=6134.4022081$, $m=0.7742237200$, and $b=-330.8452339$.
The reference values appear in Table~\ref{ba:table-reference}.

\begin{table}[htbp]
\centering\small
\begin{tabular}{@{}lll@{}}
\toprule
Quantity & Value & Units\\
\midrule
$K$ & $0.8466950351$ & $\mathrm{m^2\,s^{-2}}$\\
$z_v$ & $-0.0053932759$ & $\mathrm m$\\
$a$ & $0.1909589318$ & $\mathrm{m\,s^{-1}}$\\
$d$ & $0.0010448942$ & $\mathrm{m^2\,s^{-1}}$\\
$\Lambda$ & $5.1993222671$ & $1$\\
$\mathcal P_0$ & $1.7017148068$ & $\mathrm{m^3\,s^{-3}}$\\
$z_c$ & $0.0150057433$ & $\mathrm m$\\
$U_{H,0}$ & $6.8255509388$ & $\mathrm{m\,s^{-1}}$\\
$z_{0,w}$ & $0.0054420488$ & $\mathrm m$\\
$k_{w,0}(0)$ & $1.8887084080$ & $\mathrm{m^2\,s^{-2}}$\\
$k_{w,0}(0.1\,\mathrm m)$ & $0.9072068354$ & $\mathrm{m^2\,s^{-2}}$\\
$k_{w,0}(0.5\,\mathrm m)$ & $0.8495674195$ & $\mathrm{m^2\,s^{-2}}$\\
$\min_{[0,H]}\lambda_{\min}(R_w)$ & $0.3144828876$ & $\mathrm{m^2\,s^{-2}}$\\
\bottomrule
\end{tabular}
\caption{Budget-cutoff WJ wall-layer reference. All entries retain the
stated molecular viscosity; these are analytical reference values, not
flow-solver validation.}
\label{ba:table-reference}
\end{table}

Here $d/\nu=69.6596>2$, so the positivity and covariance propositions
apply. The gate also satisfies the stronger requirement for a tail
starting at $H/2$, since $\Lambda>1+\ln2$. At the ground,
\begin{equation}
 E_w(0)=21.81134911\,\mathrm{m^2\,s^{-3}},\qquad
 p_{A,0}(0)=58.95599772\,\mathrm{m^2\,s^{-3}}.
 \label{ba:example-local-source}
\end{equation}
The defect changes sign at $z\simeq0.0094549500\,\mathrm m$ and its
full-layer integral is zero. The corresponding pressure offset is
$(p_w(H)-p_o(H))/\rho=-0.2090605025\,\mathrm{m^2\,s^{-2}}$.

Independent split quadrature confirms the integrated production and sink,
the momentum and pressure matching, the dissipation trace, and the
mechanical-energy identity at roundoff level for this example. The first
six polynomial test moments of \eqref{ba:weak-equilibrium} have absolute
residuals below $2\times10^{-16}\,\mathrm{m^3\,s^{-3}}$ for
$\phi_n=(z/H)^n$, $n=0,\ldots,5$.
The dilogarithm expression and the integral solution for $k_{w,0}$ agree
within $3.3\times10^{-15}\,\mathrm{m^2\,s^{-2}}$ at the evaluated points.
These checks verify evaluation of the continuum formulas and do not
replace a residual test of a flow discretization.

Implementation should sample $K$ from the common matching trace, evaluate
the calibrated $D_A$ with the same positive $\nu$, and determine the cutoff
from reference production at that $K$. Use actual production in
\eqref{ba:k-transport} or \eqref{ba:curved-k}. For the flat test,
initialize \eqref{ba:k-profile}, impose zero modeled ground $k$ flux,
and couple the common matching value and conservative flux to the outer
WJ equations. The prescribed EP redistribution must not also be added
to this nonconstant-$k$ reference. Quadrature must resolve every active
viscosity, forcing, and dissipation breakpoint. The assembled reference
residual must include the actual interface, pressure, and boundary operators.

The derivation establishes continuum compatibility with the standard WJ
neutral equilibrium on flat undisturbed terrain, a positive and realizable
reference on the stated rough branch, and a conservative curved extension.
Positivity of this stationary reference is not a positivity theorem for
the evolving equation: a sink controlled by sampled $K$ can remain positive
while local $k_w$ decreases, as \eqref{ba:first-zero-failure} demonstrates.
The conditional amendment in \eqref{ba:positive-profiles} changes that
source policy and is explicitly separate from the default variant.
General-flow covariance realizability,
long-time robustness, and discrete equilibrium preservation require their
own checks. Any modification introduced for those purposes must be
declared and verified to be inactive on the stated reference or to preserve
it exactly.

\begin{thebibliography}{9}
\bibitem{WJ2000}
S.~Wallin and A.~V.~Johansson,
``An explicit algebraic Reynolds stress model for incompressible and
compressible turbulent flows,''
\emph{Journal of Fluid Mechanics} \textbf{403}, 89--132 (2000).
\href{https://doi.org/10.1017/S0022112099007004}{doi:10.1017/S0022112099007004}.

\bibitem{Baungaard2022}
M.~Baungaard, S.~Wallin, M.~P.~van der Laan and M.~Kelly,
``Wind turbine wake simulation with explicit algebraic Reynolds stress
modeling,'' \emph{Wind Energy Science} \textbf{7}, 1975--2002 (2022).
\href{https://doi.org/10.5194/wes-7-1975-2022}{doi:10.5194/wes-7-1975-2022}.

\bibitem{SCI2006}
K.~Suga, T.~J.~Craft and H.~Iacovides,
``An analytical wall-function for turbulent flows and heat transfer over
rough walls,'' \emph{International Journal of Heat and Fluid Flow}
\textbf{27}(5), 852--866 (2006).
\href{https://doi.org/10.1016/j.ijheatfluidflow.2006.03.011}{doi:10.1016/j.ijheatfluidflow.2006.03.011}.
\end{thebibliography}

\end{document}
